% =============================================================================
% TransGraph-PPI: A Hybrid Ensemble Framework for Sequence-Based
% Protein-Protein Interaction Prediction
% IEEE conference format (IEEEtran)
% =============================================================================
\documentclass[conference]{IEEEtran}
\IEEEoverridecommandlockouts

% ----------------------------------------------------------------------------
% Packages
% ----------------------------------------------------------------------------
\usepackage{cite}
\usepackage{amsmath,amssymb,amsfonts}
\usepackage{graphicx}
\usepackage{textcomp}
\usepackage{xcolor}
\usepackage{booktabs}   % cleaner table rules
\usepackage{array}

\def\BibTeX{{\rm B\kern-.05em{\sc i\kern-.025em b}\kern-.08em
    T\kern-.1667em\lower.7ex\hbox{E}\kern-.125emX}}

% ----------------------------------------------------------------------------
% Shortcuts
% ----------------------------------------------------------------------------
\newcommand{\pseq}{p_{\mathrm{seq}}}
\newcommand{\pgraph}{p_{\mathrm{graph}}}
\newcommand{\pfinal}{p_{\mathrm{final}}}

\begin{document}

% ----------------------------------------------------------------------------
% Title and authors
% ----------------------------------------------------------------------------
\title{TransGraph-PPI: A Hybrid Ensemble Framework\\for
Protein--Protein Interaction Prediction}

\author{
\begin{tabular}{ccc}
\begin{tabular}{c}
Sridevi Saralaya\\
\textit{Dept. of CSE}\\
\textit{St Joseph Engineering College}\\
Vamanjoor, India\\
sridevis@sjec.ac.in
\end{tabular}
&
\begin{tabular}{c}
Jaishma Kumari B\\
\textit{Dept. of CSE}\\
\textit{St Joseph Engineering College}\\
Vamanjoor, India\\
jaishmab@sjec.ac.in
\end{tabular}
&
\begin{tabular}{c}
Ashwini Shenoy B\\
\textit{Dept. of CSE}\\
\textit{St Joseph Engineering College}\\
Vamanjoor, India\\
23a36.ashwini@sjec.ac.in
\end{tabular}
\\[4em]
\begin{tabular}{c}
Ajay Preenal Dsouza\\
\textit{Dept. of CSE}\\
\textit{St Joseph Engineering College}\\
Vamanjoor, India\\
23a06.ajay@sjec.ac.in
\end{tabular}
&
\begin{tabular}{c}
Basil S\\
\textit{Dept. of CSE}\\
\textit{St Joseph Engineering College}\\
Vamanjoor, India\\
23a16.basil@sjec.ac.in
\end{tabular}
&
\begin{tabular}{c}
Bhavish V\\
\textit{Dept. of CSE}\\
\textit{St Joseph Engineering College}\\
Vamanjoor, India\\
23a43.bhavish@sjec.ac.in
\end{tabular}
\end{tabular}
}

\maketitle

% ----------------------------------------------------------------------------
% Abstract and keywords
% ----------------------------------------------------------------------------
\begin{abstract}
Protein--protein interactions (PPIs) are important for understanding cellular
processes, disease mechanisms, and the identification of potential drug
targets. Experimental identification of PPIs can be time-consuming and costly,
creating a need for computational methods that can predict interactions
efficiently. This paper presents TransGraph-PPI, a hybrid ensemble framework
for protein--protein interaction prediction. The framework combines protein
sequence information obtained from the ESM-2 protein language model with
graph-based information from a high-confidence STRING interaction network. A
sequence model and a graph neural network learn complementary information from
protein sequences and network structure, and their outputs are combined using
an ensemble model. SHAP-based explainability is used to identify the
contribution of different meta-features to each prediction. On the
evaluated dataset, the proposed framework achieved a ROC-AUC of 0.9708, an
accuracy of 92.13\%, and an F1-score of 0.9200. The evaluation used a
pair-disjoint split, with no overlapping protein pairs between the training,
validation, and test sets. The developed system provides an integrated
approach for predicting PPIs, explaining model predictions, and supporting
further analysis of proteins that may be relevant as therapeutic targets.
\end{abstract}

\begin{IEEEkeywords}
Protein--protein interaction, ESM-2, graph neural network, deep learning,
ensemble learning, explainable AI, SHAP, drug target prioritization
\end{IEEEkeywords}

% =============================================================================
\section{Introduction}
% =============================================================================
Protein--protein interactions (PPIs) are fundamental to biological processes
such as signal transduction, metabolism, gene regulation, and cellular
functions. Understanding PPIs is important for studying disease mechanisms,
protein functions, and identifying potential therapeutic targets. However,
experimental identification of PPIs can be expensive, time-consuming, and
difficult to scale to a large number of protein pairs. Therefore,
computational methods have been developed to predict PPIs using protein
sequences and interaction information. Recent approaches use deep learning and
protein language models (PLMs), such as ESM-2, to learn informative
representations from protein sequences, while graph neural networks (GNNs)
capture relationships between proteins in interaction networks. However,
effectively combining sequence and graph information while maintaining
interpretable predictions remains a challenge \cite{b1,b2,b3,b4}.

To address this challenge, this work proposes TransGraph-PPI, a hybrid
framework that combines ESM-2-based protein sequence representations with
graph-based learning for PPI prediction. Sequence and graph predictions are
integrated using an ensemble mechanism, while SHAP-based analysis provides
interpretation of the prediction features. The framework also supports
downstream biological analysis of predicted interactions. The main
contributions of this work are:
\begin{enumerate}
    \item integration of ESM-2 sequence representations and graph-based
          learning for PPI prediction;
    \item an ensemble prediction mechanism that combines complementary
          sequence and graph information; and
    \item an explainable analysis framework using SHAP to interpret
          predictions and support biological investigation.
\end{enumerate}

The remainder of this paper is organized as follows. Section~\ref{sec:related}
presents related work, Section~\ref{sec:design} describes the proposed
TransGraph-PPI architecture and methodology, Section~\ref{sec:impl} covers the
implementation, Section~\ref{sec:results} presents the experimental results and
discussion, and Sections~\ref{sec:future} and~\ref{sec:conclusion} discuss
future directions and conclude the paper.

% =============================================================================
\section{Literature Review}\label{sec:related}
% =============================================================================
Protein--protein interaction (PPI) prediction has progressed from traditional
machine learning with handcrafted sequence features to deep learning methods
that learn patterns directly from protein sequences. Recent approaches
increasingly use protein language models (PLMs) with convolutional and
graph-based architectures. xCAPT5 uses T5-XL-UniRef50 representations with a
multi-kernel Siamese convolutional neural network and XGBoost for binary PPI
prediction \cite{b5}. While pretrained representations reduce dependence on
handcrafted features, the method mainly captures sequence information and does
not explicitly model interaction-network topology. DeepGNHV combines ProtT5
embeddings, a Graph Attention Network (GAT), and AlphaFold2-derived structural
information for human--virus PPI prediction \cite{b6}. However, it is
specialized for human--virus interactions and requires predicted
three-dimensional structures.

Another direction examines how PLM representations can be combined for
interaction-related tasks. Alsamkary et al.\ investigated PLM architectures for
multi-chain PPI affinity prediction using ESM and ProtT5 with hierarchical
pooling and attention-based architectures \cite{b7}. Using the PPB-Affinity
dataset containing 8,207 entries, their study showed that representation and
pooling strategies can significantly affect prediction performance. However,
it focuses on binding-affinity prediction rather than binary PPI
classification and does not incorporate an interaction-network graph model,
although it highlights the importance of effective PLM representation
learning.

Graph neural networks (GNNs) have also become important in PPI-related
prediction because proteins and their relationships can be naturally
represented as graphs. ESM-GNN combines ESM-2 embeddings with a Gated Graph
Neural Network (GGNN) to predict protein--protein interaction sites
\cite{b8}. It uses sequential, K-nearest-neighbor, and predicted contact
relationships to combine sequence information with local residue
relationships. The reported results include an AUC of 0.892 and an F1-score of
0.847. However, ESM-GNN focuses on interaction-site identification rather than
binary prediction of whether two proteins interact.

PFMBench provides a broader evaluation of protein foundation models across 38
biological tasks using 17 models, including interaction, annotation,
structure, mutation, and localization tasks \cite{b9}. It compares
sequence-based and multimodal models using measures such as accuracy,
F1-score, and AUROC. The results indicate that multimodal representations can
provide advantages on several tasks and that larger models do not necessarily
guarantee better performance. However, PFMBench is a benchmarking framework
rather than an end-to-end PPI prediction architecture, providing guidance for
model evaluation and selection.

Among graph-based methods, GLTPPI is particularly relevant because it combines
PLM representations with global and local graph information \cite{b10}. It uses
ESM-2 embeddings and a Graph Attention Network with random-walk-based global
topology and a Variational Graph Autoencoder (VGAE) for local graph
structures. Evaluation on SHS27k, SHS148k, Yeast, and Arabidopsis reported
accuracies of 0.8979 and 0.9355 on SHS27k and SHS148k, respectively. GLTPPI
demonstrates that sequence information and interaction-network topology can
provide complementary information for PPI prediction. However, it does not
employ the independent sequence and graph predictors, ensemble mechanism,
interpretability, and downstream biological analysis considered in the
proposed framework.

Review studies further describe the development of deep learning for
sequence-based PPI prediction. Mewara and Lalwani surveyed deep-network
approaches, highlighting the progression from handcrafted features toward
representation learning and modern deep architectures \cite{b11}. Reim et
al.\ investigated unbiased sequence-based PPI prediction and reported that deep
learning performance can plateau under more stringent evaluation conditions
\cite{b12}, emphasizing appropriate data splits, evaluation metrics, and
generalization. Similarly, Cui et al.\ reviewed recent advances in deep
learning for PPI prediction, covering sequence, structural, PLM-based, and
graph-based approaches \cite{b13}. Their review highlights differences in
information requirements and prediction objectives, including binary PPI,
interaction-site, and binding-affinity prediction.

Overall, the reviewed studies demonstrate a progression from sequence-based
learning toward PLM-based, graph-based, and hybrid PPI prediction. PLMs such as
ESM-2, ProtT5, and T5-based models provide informative sequence
representations, while GNNs capture relationships and topology beyond
individual sequence representations. xCAPT5 demonstrates PLM-based prediction,
while DeepGNHV and ESM-GNN incorporate graph learning into protein-related
tasks \cite{b5,b6,b8}. GLTPPI further demonstrates the value of combining
sequence embeddings with global and local graph topology \cite{b10}. PFMBench
and the review studies emphasize model selection, evaluation, and
generalization \cite{b9,b11,b12,b13}.

However, existing methods often focus on a single information source or
address related tasks such as interaction-site or affinity prediction. Some
hybrid approaches require structural information, while others focus mainly on
sequence or graph topology. Interpretability is also not a central component
of most existing frameworks. These observations motivate TransGraph-PPI, which
combines ESM-2 sequence representations with graph-based learning and an
ensemble prediction mechanism. The framework integrates complementary sequence
and interaction-network information with SHAP-based analysis to identify
features contributing to predictions, providing a unified approach that
combines PLM representations, graph learning, ensemble prediction, and
explainability for PPI prediction.

% =============================================================================
\section{System Design}\label{sec:design}
% =============================================================================
TransGraph-PPI is a machine learning framework designed to predict whether two
human proteins interact with each other (protein--protein interaction, or
PPI). The framework combines information from protein sequences and the protein
interaction network. The system consists of three main components:
(1)~a sequence branch, which predicts interactions from protein sequence
information; (2)~a graph branch, which predicts interactions using the protein
interaction network; and (3)~an ensemble branch, which combines the
predictions using XGBoost. SHAP is then used to explain the final predictions.

\subsection{Overall Architecture}\label{AA}
Protein sequences are first processed using ESM-2
(\texttt{facebook/esm2\_t12\_35M\_UR50D}), which produces a 480-dimensional
embedding for each protein. These embeddings are provided to the sequence and
graph branches. The sequence branch uses a residual multi-layer perceptron
(MLP) to learn interaction patterns from the two protein embeddings. The graph
branch uses GraphSAGE to capture information from neighboring proteins in the
interaction network. The two branches produce separate interaction
probabilities, $\pseq$ and $\pgraph$.

Instead of directly averaging these probabilities, TransGraph-PPI constructs an
8-dimensional feature vector containing the two prediction scores, their
confidence values, their difference, their agreement, and biological
information. An XGBoost model uses these features to produce the final
interaction probability.

\subsection{Sequence-Based Prediction}\label{sec:seq}
Since the interaction between proteins $A$ and $B$ is the same as the
interaction between $B$ and $A$, the model should not depend on the order of
the proteins. Therefore, the two 480-dimensional ESM-2 embeddings, $e_A$ and
$e_B$, are combined using four symmetric operations:
\begin{equation}
f_{\mathrm{sym}} =
\bigl[
(e_A + e_B)
\parallel
(e_A \odot e_B)
\parallel
|e_A - e_B|
\parallel
\max(e_A,e_B)
\bigr]
\label{eq:sym}
\end{equation}
where $\parallel$ represents concatenation and $\odot$ represents element-wise
multiplication. The four operations produce a 1920-dimensional feature vector.

This vector is passed through a residual MLP containing skip connections,
normalization, and dropout. The model produces a probability $\pseq$ between 0
and 1, representing the likelihood that the two proteins interact based on
their sequence information.

\subsection{Graph-Based Prediction}\label{sec:graph}
Proteins form a complex biological network, so network information can provide
additional evidence for PPI prediction. The graph branch represents proteins as
nodes and known interactions from the STRING database as edges.

Each protein node is represented using a 483-dimensional feature vector,
consisting of the 480-dimensional ESM-2 embedding and three network features:
degree centrality, clustering coefficient, and PageRank.

GraphSAGE updates each protein representation by aggregating information from
its neighboring proteins:
\begin{equation}
h_i^{(l)} =
\operatorname{ReLU}\!\left(
W_1 h_i^{(l-1)}
+
W_2 \cdot
\operatorname*{mean}_{j \in \mathcal{N}(i)} h_j^{(l-1)}
\right)
\label{eq:sage}
\end{equation}
where $\mathcal{N}(i)$ represents the neighboring proteins of protein $i$, and
$W_1$ and $W_2$ are learnable weights.

The learned representations of two candidate proteins are then combined using
several pairwise operations and passed through a classifier to obtain the
graph-based interaction probability $\pgraph$.

\subsection{XGBoost Stacking Ensemble}\label{sec:stack}
The sequence and graph models capture different types of information. The
sequence model focuses on protein sequence patterns, while the graph model
captures information from the interaction network. To combine these
predictions, TransGraph-PPI uses an XGBoost stacking model.

To reduce data leakage, 5-fold stratified out-of-fold (OOF) predictions are
used to train the ensemble model. For each protein pair, eight meta-features
are calculated:
\begin{equation}
\begin{split}
X_{\mathrm{meta}} = [\,&\pseq,\ \pgraph,\ \mathrm{conf}_{\mathrm{seq}},\
\mathrm{conf}_{\mathrm{graph}},\\
&\mathrm{diff},\ \mathrm{max}_{\mathrm{conf}},\
\mathrm{consensus},\ \mathrm{bio}_{\mathrm{score}}\,]
\end{split}
\label{eq:meta}
\end{equation}
The eight features are defined as follows:
\begin{itemize}
    \item $\pseq$: probability predicted by the sequence model.
    \item $\pgraph$: probability predicted by the graph model.
    \item $\mathrm{conf}_{\mathrm{seq}} = |\pseq - 0.5|$: confidence of the
          sequence model.
    \item $\mathrm{conf}_{\mathrm{graph}} = |\pgraph - 0.5|$: confidence of
          the graph model.
    \item $\mathrm{diff} = |\pseq - \pgraph|$: difference between the two
          predictions.
    \item $\mathrm{max}_{\mathrm{conf}} =
          \max(\mathrm{conf}_{\mathrm{seq}},\mathrm{conf}_{\mathrm{graph}})$:
          higher confidence of the two models.
    \item $\mathrm{consensus} = \pseq \times \pgraph$: agreement between the
          two predictions.
    \item $\mathrm{bio}_{\mathrm{score}}$: biological compatibility score
          based on subcellular co-localization.
\end{itemize}
The XGBoost model uses these meta-features to produce the final interaction
probability:
\begin{equation}
\pfinal = \operatorname{XGBoost}(X_{\mathrm{meta}})
\label{eq:final}
\end{equation}

\subsection{Model Explainability}\label{sec:shap}
SHAP (SHapley Additive exPlanations) is applied to the trained XGBoost model to
make the ensemble predictions easier to interpret. SHAP measures how much each
of the eight meta-features contributes to increasing or decreasing the final
prediction score. This helps identify whether a prediction is mainly supported
by sequence information, network information, or agreement between the two
models.

% =============================================================================
\section{Implementation}\label{sec:impl}
% =============================================================================
The TransGraph-PPI framework was implemented using Python and deep learning
libraries for sequence representation, graph learning, ensemble prediction, and
explainability. The implementation follows the architecture described in
Section~\ref{sec:design} and includes dataset preparation, sequence
representation, model training, graph construction, ensemble learning, and
SHAP-based analysis.

\subsection{Dataset Preparation}
The dataset contains 201,712 human protein pairs, including 100,856 positive
and 100,856 negative samples. Positive interactions were obtained from the
STRING database using a minimum confidence score of 900. Negative pairs were
sampled while excluding known positive interactions to avoid incorrect negative
labels.

The dataset was divided using a stratified 80:10:10 split, resulting in 161,369
training pairs, 20,171 validation pairs, and 20,172 test pairs. The validation
set was used for threshold selection, while the test set was kept separate for
final evaluation.

\subsection{Protein Sequence Representation}
Each protein sequence was represented using the pretrained ESM-2 model
\texttt{facebook/esm2\_t12\_35M\_UR50D}. The model contains approximately 35
million parameters and produces a 480-dimensional mean-pooled representation
for each protein.

These embeddings were used as inputs to both the sequence and graph branches,
allowing sequence information to be combined with interaction-network
information without manually designing extensive sequence features.

\subsection{Sequence Model Implementation}
The sequence branch was implemented using a residual MLP. The symmetric pair
representation described in Section~\ref{sec:seq} produces a 1920-dimensional
input. This representation is first projected to 1024 dimensions, followed by
two residual blocks that progressively reduce the representation to 512
dimensions before the final classification layers.

The model was trained using AdamW with a learning rate of $10^{-3}$ and weight
decay of $10^{-3}$. A cosine annealing learning-rate scheduler was used,
together with balanced focal loss to emphasize difficult predictions. The batch
size was 64, with a maximum of 30 epochs and early stopping.

\subsection{Graph Model Implementation}
For the graph branch, proteins were represented as nodes and high-confidence
interactions as edges. Each node contained the 480-dimensional ESM-2 embedding
together with degree centrality, clustering coefficient, and PageRank,
resulting in 483 node features.

A two-layer GraphSAGE model with a hidden dimension of 256 was used for
neighborhood aggregation. Dropout and normalization were applied to reduce
overfitting. For candidate protein pairs, the learned node representations were
combined using pairwise operations and a bilinear interaction layer before
classification.

The graph model was trained using AdamW with a learning rate of $10^{-3}$,
weight decay of $10^{-4}$, and a cosine annealing scheduler. Focal loss was
used as the training objective. Graph edges were processed in chunks during
link prediction to efficiently handle the large interaction graph.

\subsection{Ensemble Implementation}
The outputs of the sequence and graph models were combined using an XGBoost
classifier. To prevent target leakage, five-fold stratified OOF predictions
were generated from the base models before training the meta-learner.

For each protein pair, the two prediction probabilities and the six additional
features described in Section~\ref{sec:stack} formed an eight-dimensional input
vector for XGBoost.

The XGBoost model used 500 estimators, a maximum tree depth of 7, a learning
rate of 0.01, and a subsampling ratio of 0.8. Feature subsampling was also set
to 0.8, with $L_1$ regularization applied to reduce overfitting.

\subsection{Explainability}
After training, SHAP was used to analyze the contribution of the eight
meta-features to the final prediction. This analysis provides insight into how
sequence predictions, graph predictions, model confidence, prediction
agreement, and biological compatibility influence the ensemble output.

% =============================================================================
\section{Results and Discussion}\label{sec:results}
% =============================================================================
The performance of TransGraph-PPI was evaluated on a held-out test set of
20,172 protein pairs, with no overlap between training and test pairs. The
sequence model, graph model, and proposed XGBoost ensemble were compared with a
Random Forest baseline using accuracy, precision, recall, F1-score, ROC-AUC,
and PR-AUC.

\subsection{Performance Comparison}

\begin{table}[htbp]
\caption{Performance Comparison of the Evaluated Models}
\label{tab:performance}
\centering
\resizebox{\columnwidth}{!}{%
\begin{tabular}{lcccccc}
\toprule
\textbf{Model} & \textbf{Accuracy} & \textbf{Precision} & \textbf{Recall} &
\textbf{F1-Score} & \textbf{ROC-AUC} & \textbf{PR-AUC} \\
\midrule
Random Forest    & 82.34\%          & 81.14\%          & 84.28\% & 82.68\%          & 0.9065          & 0.9124 \\
ESM-MLP          & 87.11\%          & 86.00\%          & 88.65\% & 87.30\%          & 0.9438          & 0.9474 \\
GraphSAGE        & 90.42\%          & 92.01\%          & 88.52\% & 90.23\%          & 0.9487          & 0.9616 \\
XGBoost Ensemble & \textbf{92.13\%} & \textbf{93.59\%} & 90.46\% & \textbf{92.00\%} & \textbf{0.9708} & \textbf{0.9759} \\
\bottomrule
\end{tabular}%
}
\end{table}

The Random Forest baseline achieved 82.34\% accuracy and an 82.68\% F1-score,
while the ESM-MLP improved these values to 87.11\% and 87.30\%, respectively,
demonstrating the usefulness of pretrained ESM-2 representations. GraphSAGE
achieved 90.42\% accuracy and a 90.23\% F1-score, with a PR-AUC of 0.9616,
indicating the contribution of interaction-network information.

The XGBoost stacking ensemble achieved the highest performance, with 92.13\%
accuracy, 92.00\% F1-score, 0.9708 ROC-AUC, and 0.9759 PR-AUC. Its precision of
93.59\% indicates that most interactions classified as positive were correct.

\subsection{Effect of Ensemble Learning}
Combining the sequence and graph branches improves overall prediction
performance. The ESM-MLP captures sequence-level information, whereas
GraphSAGE incorporates interaction-network information, allowing the ensemble
to exploit complementary signals.

Compared with ESM-MLP, the ensemble improves accuracy from 87.11\% to 92.13\%
and F1-score from 87.30\% to 92.00\%. Compared with GraphSAGE, accuracy
increases from 90.42\% to 92.13\% and ROC-AUC from 0.9487 to 0.9708. These
results support the use of stacking to combine the two prediction branches. The
ensemble also achieves a PR-AUC of 0.9759, indicating good precision-recall
performance across classification thresholds. Furthermore, training the XGBoost
meta-learner using OOF predictions reduces the risk of target leakage during
stacking.

\subsection{Discussion}
The results demonstrate that both sequence and interaction-network information
contribute to PPI prediction. The improvement from ESM-MLP to GraphSAGE
suggests that network topology provides information beyond individual sequence
representations, while the further improvement from stacking indicates that the
two branches provide complementary signals.

The ensemble achieves higher precision than the individual models, which is
relevant for PPI prediction because false-positive interactions may affect
downstream biological analysis. By combining predictions and confidence-related
features, the XGBoost meta-learner can make the final prediction using multiple
sources of evidence.

SHAP analysis provides an additional interpretability layer by identifying the
contribution of individual meta-features to XGBoost predictions. Overall, the
results demonstrate that integrating ESM-2 sequence representations,
GraphSAGE-based network learning, and XGBoost stacking provides an effective
framework for PPI prediction while also supporting interpretation of the
resulting predictions.

% =============================================================================
\section{Future Work}\label{sec:future}
% =============================================================================
Future work will focus on improving the generalization and biological
usefulness of TransGraph-PPI. The framework can be extended by incorporating
additional protein sequence and structural information, larger and more diverse
interaction datasets, and more advanced graph learning methods. Further
experiments can also evaluate the model on unseen proteins and independent
datasets to assess its generalization to different biological settings. The
explainability component can be extended to provide more detailed biological
interpretation of important interaction features. These improvements may help
make the framework more useful for large-scale PPI analysis and downstream
drug-target research.

% =============================================================================
\section{Conclusion}\label{sec:conclusion}
% =============================================================================
TransGraph-PPI presents a hybrid framework for predicting protein--protein
interactions by combining protein sequence information with
interaction-network information. The framework uses ESM-2 embeddings to
represent protein sequences, a residual MLP for sequence-based prediction, and
GraphSAGE for learning information from the protein interaction network. The
predictions from both branches are combined using an XGBoost stacking model
with eight meta-features, while SHAP provides an interpretable view of the
final predictions.

On the held-out test set, the proposed ensemble achieved 92.13\% accuracy,
92.00\% F1-score, 0.9708 ROC-AUC, and 0.9759 PR-AUC, outperforming the
individual sequence and graph branches. These results show that combining
complementary sequence and network information can improve PPI prediction
compared with using either source independently. The use of SHAP further
supports interpretation of the ensemble predictions. Overall, TransGraph-PPI
provides an integrated approach for PPI prediction and explainable analysis,
with potential applications in studying protein interactions and supporting
downstream biological research.

% =============================================================================
% References
% =============================================================================
\begin{thebibliography}{00}

\bibitem{b1}
F. Charih, J. R. Green, and K. K. Biggar,
``Sequence-based protein--protein interaction prediction and its applications
in drug discovery,''
\textit{Cells}, vol.~14, no.~18, Art.~no.~1449, 2025,
doi: 10.3390/cells14181449.

\bibitem{b2}
N. Alkhateeb and M. Awad,
``Advances in protein--protein interaction prediction: A deep learning
perspective,'' 2023.

\bibitem{b3}
W. Hu and M. Ohue,
``SpatialPPIv2: Enhancing protein--protein interaction prediction through graph
neural networks with protein language models,''
\textit{Comput.\ Struct.\ Biotechnol.\ J.}, vol.~27, pp.~508--518, 2025,
doi: 10.1016/j.csbj.2025.01.022.

\bibitem{b4}
F. Soleymani, E. Paquet, H. L. Viktor, W. Michalowski, and D. Spinello,
``ProtInteract: A deep learning framework for predicting protein--protein
interactions,''
\textit{Comput.\ Struct.\ Biotechnol.\ J.}, vol.~21, pp.~1324--1348, 2023,
doi: 10.1016/j.csbj.2023.01.028.

\bibitem{b5}
T. H. Dang and T. A. Vu,
``xCAPT5: Protein--protein interaction prediction using deep and wide
multi-kernel pooling convolutional neural networks with protein language
model,''
\textit{BMC Bioinformatics}, vol.~25, Art.~no.~106, 2024,
doi: 10.1186/s12859-024-05725-6.

\bibitem{b6}
L. Jiang, X. Yang, X. Guo, D. Li, J. Li, S. Wuchty, W. Shi, and Z. Zhang,
``Graph neural network integrated with pretrained protein language model for
predicting human-virus protein--protein interactions,''
\textit{Briefings in Bioinformatics}, vol.~26, no.~5, Art.~no.~bbaf461, 2025,
doi: 10.1093/bib/bbaf461.

\bibitem{b7}
H. Alsamkary, M. Elshaffei, M. Soudy, S. Ossman, A. Amr, N. A. Abdelsalam,
M. Elkerdawy, and A. Elnaggar,
``Beyond simple concatenation: Fairly assessing {PLM} architectures for
multi-chain protein-protein interactions prediction,''
\textit{arXiv preprint arXiv:2505.20036}, 2025.

\bibitem{b8}
M. K. P. and A. Geetha,
``Predicting human protein--protein interaction sites by integrating {ESM-2}
sequence embeddings and a gated graph neural network,'' 2024.

\bibitem{b9}
Z. Gao, H. Wang, C. Tan, C. Xu, M. Liu, B. Hu, L. Chao, X. Zhang, and
S. Z. Li,
``{PFMBench}: Protein foundation model benchmark,''
\textit{arXiv preprint arXiv:2506.14796}, 2025.

\bibitem{b10}
J. Liu and T. Zhang,
``{ESM2}-driven protein-protein interaction prediction with global-local graph
topology learning,''
in \textit{Proc.\ IEEE Int.\ Joint Conf.\ Neural Networks (IJCNN)}, 2025,
pp.~1--8, doi: 10.1109/IJCNN64981.2025.11228802.

\bibitem{b11}
B. Mewara and S. Lalwani,
``A survey on deep networks approaches in prediction of sequence-based
protein--protein interactions,''
\textit{SN Comput.\ Sci.}, vol.~3, no.~4, Art.~no.~298, 2022,
doi: 10.1007/s42979-022-01197-8.

\bibitem{b12}
T. Reim, A. Hartebrodt, D. B. Blumenthal, J. Bernett, and M. List,
``Deep learning models for unbiased sequence-based {PPI} prediction plateau at
an accuracy of 0.65,''
\textit{Bioinformatics}, vol.~41, no.~1, pp.~i590--i598, 2025,
doi: 10.1093/bioinformatics/btaf192.

\bibitem{b13}
J. Cui, S. Yang, L. Yi, Q. Xi, D. Yang, and Y. Zuo,
``Recent advances in deep learning for protein-protein interaction: A
review,''
\textit{BioData Mining}, vol.~18, Art.~no.~43, 2025,
doi: 10.1186/s13040-025-00457-6.

\end{thebibliography}

\end{document}